\documentclass[11pt,a4paper]{article}
% Common preamble for the hint sheets (\input from each file)
\usepackage[utf8]{inputenc}
\usepackage[T1]{fontenc}
\usepackage{lmodern}
\usepackage{amsmath,amssymb,amsthm,mathtools}
\usepackage[a4paper,margin=2.3cm]{geometry}
\usepackage{enumitem}
\usepackage{xcolor}
\usepackage[hidelinks]{hyperref}
\renewcommand{\leq}{\leqslant}
\renewcommand{\le}{\leqslant}
\renewcommand{\geq}{\geqslant}
\renewcommand{\ge}{\geqslant}
\setlength{\parindent}{0pt}
\setlength{\parskip}{0.35em}
% \hintstitle{Sheet title}{Class work / Homework}
\newcommand{\hintstitle}[2]{\begin{center}{\LARGE\bfseries Hints: #1\par}\vspace{0.5em}{\large #2\par}\end{center}
\noindent\textit{The hints to each problem are cumulative. Read one, try again, and open the next only if you are still stuck.}\par\medskip\hrule\medskip}
\newcommand{\problem}[1]{\par\bigskip{\large\bfseries Problem #1}\par\nopagebreak\smallskip}
\newcommand{\subpart}[1]{\par\smallskip\textbf{(#1)}\par\nopagebreak}
\newcommand{\hint}[1]{\par\smallskip\noindent\textcolor{black!70}{\textbf{Hint #1.}}\enspace}
\newcommand{\note}[1]{\par\smallskip\noindent\textit{Note.} #1\par}
\providecommand{\lt}{<}
\providecommand{\gt}{>}
\providecommand{\tightlist}{\setlength{\itemsep}{0pt}\setlength{\parskip}{0pt}}
\emergencystretch=3em

\begin{document}
\hintstitle{Methods of proof -- 1}{Homework}

\problem{1}
\note{``Rationally independent $\alpha_i$'' is not enough for Kronecker's theorem. For example, $\sqrt2$ and $1+\sqrt2$ are linearly independent over $\mathbb Q$, but the points $(\{q\sqrt2\},\{q(1+\sqrt2)\})$ all lie on the diagonal. The correct hypothesis is that $1,\alpha_1,\dots,\alpha_k$ are linearly independent over $\mathbb Q$, and the hints use it.}
\hint{1} Cut the unit cube $[0,1)^k$ into $N^k$ small cubes of side $1/N$, and place the $N^k+1$ points $(\{q\alpha_1\},\dots,\{q\alpha_k\})$, $q=0,1,\dots,N^k$, into them.
\hint{2} Two points with indices $a<b$ fall into the same small cube. Then $q=b-a\le N^k$ satisfies $\|q\alpha_i\|<1/N$ for every $i$.
\hint{3} For density, recurrence near $0$ is not enough when $k\ge 2$: multiples of one small vector only fill a line. Use Weyl's criterion instead. For an integer vector $h\neq0$, the number $h\cdot\alpha$ is not an integer (this is exactly where the independence of $1,\alpha_1,\dots,\alpha_k$ enters). So $\frac1Q\sum_{q=1}^Q e^{2\pi i q\,h\cdot\alpha}$ is a geometric sum bounded by $\frac{2}{Q|1-e^{2\pi i h\cdot\alpha}|}\to0$.
\hint{4} Hence $\frac1Q\sum_{q\le Q}P(q\alpha)\to\int_{[0,1]^k}P$ for every trigonometric polynomial $P$. Given a small open box $B$, take a continuous $0\le\varphi\le\mathbf 1_B$ with $\int\varphi>0$ and approximate it uniformly by a trigonometric polynomial (Fej\'er or Stone--Weierstrass). The averages of $\varphi(q\alpha)$ have a positive limit, so some $q\alpha$ lies in $B$ modulo $\mathbb Z^k$.

\problem{2}
\note{The optimal constant $d_n$ (the largest possible minimum distance of $n$ points in a unit square) is not known in closed form for general $n$. It is the classical problem of packing equal circles in a square, solved only for small $n$. The hints cover the case $n=5$, the order of magnitude $d_n\asymp n^{-1/2}$, and the true asymptotics.}
\hint{1} For five points, cut the square into four squares of side $\tfrac12$. Two points lie in the same small closed square, so their distance is at most its diagonal, $\sqrt2/2$. The bound is attained by the four corners together with the centre.
\hint{2} Upper bound for $n$ points: if $m^2<n$, cut the square into $m^2$ squares of side $1/m$. This gives $d_n\le\sqrt2/m$, so $d_n\le\sqrt2/(\lceil\sqrt n\rceil-1)\sim\sqrt{2/n}$. Lower bound: choose $n$ points of the $k\times k$ grid with spacing $1/(k-1)$, where $k=\lceil\sqrt n\rceil$. This gives $d_n\ge 1/(\lceil\sqrt n\rceil-1)\ge1/\sqrt n$.
\hint{3} A better upper bound comes from area instead of boxes. The disks of radius $d_n/2$ around the points are disjoint and lie in the square enlarged by $d_n/2$ on each side, so $n\pi d_n^2/4\le(1+d_n)^2$. This gives $d_n\le(1+o(1))\,2/\sqrt{\pi n}$.
\hint{4} The exact asymptotics is $d_n\sim\sqrt{2/(\sqrt3\,n)}$. The lower bound comes from a hexagonal (triangular lattice) arrangement. The upper bound needs Thue's theorem that the density of a packing of equal disks is at most $\pi/(2\sqrt3)$. Boundary effects are $O(\sqrt n)$ disks out of $n$.

\problem{3}
\hint{1} $2^n$ begins with the digit string $2026$ if and only if $2026\cdot10^j\le 2^n<2027\cdot10^j$ for some $j$. Take $\log_{10}$: this says that the fractional part $\{n\log_{10}2\}$ lies in $[\log_{10}2.026,\ \log_{10}2.027)$.
\hint{2} The number $\log_{10}2$ is irrational: $2^a=10^b$ with $a\geq 1$ is impossible, because $5\mid 10^b$ but $5\nmid 2^a$. So the sequence $\{n\log_{10}2\}$ is dense in $[0,1]$ (pigeonhole, as in the class work), and it meets every interval of positive length.
\hint{3} The same argument works for any digit string $D$ whose first digit is non-zero: use the interval $[\log_{10}(D/10^{r}),\log_{10}((D+1)/10^{r}))$, where $r+1$ is the number of digits of $D$. So exactly the strings that do not start with $0$ occur.

\problem{4}
\hint{1} Density theorem: for measurable $A\subseteq\mathbb R$, almost every $x\in A$ satisfies $m(A\cap I)/|I|\to1$ as intervals $I\ni x$ shrink to $x$. It suffices to show that for each $\varepsilon>0$ the set $E$ of points of $A$ where $\liminf m(A\cap I)/|I|<1-\varepsilon$ has measure $0$.
\hint{2} Suppose $m^*(E)>0$ (assume $A$ is bounded). Take an open $U\supseteq A$ with $m(U)<m(A)+\eta$. Each $x\in E$ lies in arbitrarily small intervals $I\subseteq U$ with $m(I\setminus A)>\varepsilon|I|$. By the Vitali covering lemma, choose disjoint such intervals that cover $E$ up to a null set. These intervals have total length at least $m^*(E)$, so $\varepsilon\, m^*(E)\le\sum m(I\setminus A)\le m(U\setminus A)<\eta$, which is a contradiction for small $\eta$.
\hint{3} Steinhaus: take a density point $x$ of $A$ and an interval $I$ around $x$ with $m(A\cap I)>\tfrac34|I|$. For $|t|<|I|/2$, the sets $A\cap I$ and $(A\cap I)+t$ both lie in $I\cup(I+t)$, which has length $|I|+|t|<\tfrac32|I|$. Their measures add up to more than $\tfrac32|I|$, so they intersect, and $t\in A-A$.
\hint{4} Both proofs extend to $\mathbb R^n$: the density theorem holds there with balls or cubes, by the Vitali lemma in $\mathbb R^n$. The regularity proof from the class work ($K\subseteq A$ compact, $U\supseteq K$ open with $m(U)<2m(K)$, and $K+t\subseteq U$ for small $t$) needs no covering lemma at all. It extends further, to every locally compact group with Haar measure, where density theorems can fail.

\problem{5}
\note{(1) For odd $n$ the first claim is false: $\{\frac{n+1}2,\dots,n\}$ has $\frac{n+1}2>\frac n2$ elements, and none of them divides another (e.g.\ $\{2,3\}$ for $n=3$). The correct hypothesis is $|A|>\lceil n/2\rceil$, which is the same as $|A|>n/2$ for even $n$. (2) As stated, the second claim holds as soon as $|A|\ge2$, and the bound $2n/3$ plays no role. The hints prove both claims in these forms.}
\hint{1} Use the odd parts as boxes, as in the class work. Writing $a=2^s u$ with $u$ odd, there are $\lceil n/2\rceil$ possible values of $u\le n$. Two elements with the same odd part divide one another.
\hint{2} For $A+A$: take any two elements $a<b$ of $A$ and look at $a+a$, $a+b$, $b+b$.

\problem{6}
\note{The first claim is false: the $n+1$ vectors from the centre of a regular simplex to its vertices make pairwise obtuse angles. The correct statement is that any $n+2$ vectors in $\mathbb R^n$ contain two with angle at most $\pi/2$, i.e.\ at most $n+1$ vectors can be pairwise obtuse. This gives the second claim for $n\ge2$, since $2^n\ge n+2$. For $n=1$ the vectors $1,-1$ are a counterexample. The hints prove the bound $n+1$.}
\hint{1} Pairwise obtuse means $\langle v_i,v_j\rangle<0$ for all $i\neq j$. Induct on $n$. In $\mathbb R^1$ at most two non-zero numbers can have pairwise negative products.
\hint{2} Given pairwise obtuse $v_0,\dots,v_m$ in $\mathbb R^n$, project $v_1,\dots,v_m$ onto $v_0^\perp$: write $v_i=w_i+c_iv_0$, where $c_i=\langle v_i,v_0\rangle/|v_0|^2<0$.
\hint{3} Then $\langle w_i,w_j\rangle=\langle v_i,v_j\rangle-c_ic_j|v_0|^2<0$. So $w_1,\dots,w_m$ are pairwise obtuse (in particular non-zero) vectors in an $(n-1)$-dimensional space. By induction $m\le n$, so there are at most $n+1$ vectors. The simplex shows the bound is sharp.

\problem{7}
\note{The inequality $X\ge m-\lambda\sigma$ only bounds the lower tail and cannot by itself push a random cut above its mean $m/2$. The hints prove the inequality and then get the excess $c\sqrt m$ from a lower bound on $\mathrm E|X-\mathrm EX|$, which is again a moment argument.}
\hint{1} The inequality is Chebyshev: $P(X<m-\lambda\sigma)\le P(|X-m|\ge\lambda\sigma)\le \operatorname{Var}X/(\lambda\sigma)^2\le1/\lambda^2$.
\hint{2} Put independent random signs $\varepsilon_v=\pm1$ on the vertices and let $X$ be the number of edges whose endpoints get different signs. Then $X=\frac m2-\frac12Y$ with $Y=\sum_{uv\in E}\varepsilon_u\varepsilon_v$, $\mathrm EY=0$ and $\mathrm EY^2=m$, because the products $\varepsilon_u\varepsilon_v$ for distinct edges are orthogonal.
\hint{3} Bound $\mathrm EY^4$. Expanding, only ordered 4-tuples of edges in which every vertex appears an even number of times contribute $1$. These are the tuples of the form $e,e,f,f$ (at most $3m^2$ of them) and the 4-cycles. A pair of disjoint opposite edges lies in at most two 4-cycles, so the 4-cycles contribute $O(m^2)$. Hence $\mathrm EY^4\le Cm^2$.
\hint{4} By H\"older, $(\mathrm EY^2)^3\le(\mathrm E|Y|)^2\,\mathrm EY^4$, so $\mathrm E|Y|\ge m^{3/2}/(\sqrt{C}\,m)=c'\sqrt m$. Since $\mathrm EY=0$, the negative part satisfies $\mathrm EY^-=\tfrac12\mathrm E|Y|\geq \tfrac{c'}{2}\sqrt m$, so $Y\le-\tfrac{c'}2\sqrt m$ for some choice of signs. The cut defined by those signs has at least $\frac m2+\frac{c'}4\sqrt m$ edges.

\problem{8}
\hint{1} Infinite pigeonhole: if each of the finitely many colour classes were finite, their union would be finite.
\hint{2} Infinite Ramsey for a 2-colouring (or $r$-colouring) of the pairs of $\mathbb N$: pick $x_1$. By pigeonhole, infinitely many $y$ have the same colour $c_1$ on $\{x_1,y\}$; call this infinite set $S_1$. Pick $x_2\in S_1$ and repeat inside $S_1$. This gives $x_1,x_2,\dots$ such that the colour of $\{x_i,x_j\}$ with $i<j$ depends only on $i$ (it equals $c_i$).
\hint{3} Apply pigeonhole once more to the colours $c_i$: infinitely many $i$ share a colour, and those $x_i$ form a monochromatic infinite set.
\hint{4} Finite Ramsey from infinite Ramsey: suppose that for some $k$, every $N$ has a colouring of $K_N$ with no monochromatic $K_k$. These colourings form a finitely branching tree (restrict a colouring of $K_{N+1}$ to $K_N$). By K\"onig's lemma (compactness of $\{0,1\}^{\binom{\mathbb N}{2}}$) it has an infinite branch, i.e.\ a colouring of all pairs of $\mathbb N$ with no monochromatic $K_k$. This contradicts the infinite theorem.

\problem{9}
\hint{1} Case $k=1$ (and the main lemma). Let $\delta>0$ be the upper density and $L+1>1/\delta$. Show that among any integers $n_0<\dots<n_L$, some difference $n_i-n_j$ lies in $A-A$. Consider the translates $A+n_i$ restricted to $[1,N]$, for $N$ along a sequence with $|A\cap[1,N]|\ge(\delta-\varepsilon)N$.
\hint{2} These $L+1$ translates lie in $[1,N+n_L]$ (shift so that $n_0\geq 0$) and have total size at least $(L+1)(\delta-\varepsilon)N>N+n_L$ for large $N$. So two of them meet: $a+n_i=a'+n_j$, i.e.\ $n_j-n_i\in A-A$, and it is non-zero.
\hint{3} General $k$: suppose that for every $d\ge1$ some $j(d)\in\{1,\dots,k\}$ has $j(d)\,d\notin A-A$. Colour each pair $\{x<y\}$ of positive integers with the colour $j(y-x)$. By (finite) Ramsey's theorem with $k$ colours, there are $L+1$ integers $x_0<\dots<x_L$ all of whose pairs have the same colour $j$.
\hint{4} Then no difference of the $L+1$ integers $jx_0,\dots,jx_L$ lies in $A-A$, which contradicts the lemma of Hints 1--2. So some $d$ has $d,2d,\dots,kd\in A-A$. The general case needs exactly this extra Ramsey step.

\problem{10}
\hint{1} Erd\H{o}s--Szekeres: label $a_i$ with the pair $(\ell_i,r_i)$, the lengths of the longest increasing and longest decreasing subsequences ending at $a_i$. Distinct indices get distinct pairs: if $i<j$, then $a_j$ extends one of the two sequences ending at $a_i$. With all $\ell_i,r_i\le n$ there are only $n^2$ possible pairs.
\hint{2} Sharpness: take $n$ blocks, each block a decreasing run of $n$ numbers, with every block lying entirely above the previous one, e.g.\ $n,n-1,\dots,1,\ 2n,\dots,n+1,\ \dots$. An increasing subsequence takes at most one term from each block, and a decreasing one stays inside a single block.
\hint{3} Random permutation of $\{1,\dots,N\}$: by Erd\H{o}s--Szekeres, $\max(\mathrm{LIS},\mathrm{LDS})\ge\sqrt N$ always. Reversing the permutation swaps LIS and LDS, so $P(\mathrm{LIS}\ge\sqrt N)\ge\tfrac12$.
\hint{4} For the matching upper bound (first moment): the expected number of increasing subsequences of length $\ell$ is $\binom N\ell/\ell!\le(e^2N/\ell^2)^\ell$, which tends to $0$ for $\ell\ge(1+\varepsilon)e\sqrt N$. So the LIS of a random permutation is of order $\sqrt N$ with high probability.

\end{document}
